We will now define the classical theory of 
$K_1(R)$ and $K_2(R)$ groups for a unital ring $R$.
As with $K_0(R)$, their definitions require only elementary linear algebra and group theory.

For reference, [Bass] and [Milnor] are excellent sources on classical $K$-theory.
%However, in the upcoming Chapter 4, we will define groups $K_n(R)$ for $n \geq 0$
%using homotopy thoeretic machinery.
%We will then need to show that these definitions match the theory for $n = 0,1,2$. 

\section{The Group $K_1$ of a ring}

Let $R$ be a unital ring. Consider the groups $GL_n(R)$ for $n \geq 0$.
One may inject each $GL_n(R)$ into $GL_{n+1}(R)$ via the embedding
\begin{align*}
  GL_n(R) &\hookrightarrow GL_{n+1}(R) \\
  g &\mapsto \pmat g 0 0 1
\end{align*}

This gives us a filtered system of groups
\[
R^\times \cong GL_1(R) \subset GL_2(R)\subset \cdots \subset GL_n(R) \subset GL_{n+1}(R) \subset \cdots 
\]

\begin{definition}
Taking the colimit of this filtered system gives us the \emph{stable general linear group},
which we denote $GL(R)$.
\end{definition}

\begin{definition}[Elementary matrices and subgroup]
Suppose $i \neq j$, the the elementary matrix $e_{ij}$ is the matrix in $GL_n(R)$ given by  
$e_{ij}(r) = I_n + r E_{ij}$, that is, identity along the diagonal and $r$ in the $(i,j)$-th spot.

The subgroup of $GL_n(R)$ generated by $e_{ij}(r)$ is called the \emph{elementary subgroup},
denoted $E_n(R)$.
\end{definition}

Similar to $GL(R)$, we have the embeddings $E_n(R) \hookrightarrow E_{n+1}(R)$ given by
\begin{align*}
E_n(R) &\hookrightarrow E_{n+1}(R) \\
e &\mapsto \pmat e 0 0 1
\end{align*}

This gives us another filtered system of groups
\[
1 \cong E_1(R) \subset E_2(R)\subset \cdots \subset E_n(R) \subset E_{n+1}(R) \subset \cdots 
\]


\begin{definition}
The colimit of this filtered system gives us the \emph{stable elementary subgroup},
denoted $E(R)$, which may be identified with a subgroup of the stable group $GL(R)$.
\end{definition}

We prove an elementary property about elementary matrices. 

\begin{lemma}
Suppose $g$ is an element in $GL_n(R)$. Then
$\psmat g 0 0 {g^{-1}}$ is in $E_{2n}(R)$.
\end{lemma}
\begin{proof}
Note that matrices of the form $\psmat 1 g 0 1$ are in $E_{2n}(R)$.
Consequently, we have the factorization
\[
\pmat g 0 0 {g^{-1}} = \pmat 1 g 0 1 \pmat 1 0 {g^{-1}} 1 \pmat 1 g 0 1 \pmat 0 {-1} 1 0
\]
so the given element is in $E_{2n}(R)$.
\end{proof}

\begin{proposition}
$E(R)$ is a normal subgroup.
\end{proposition}
\begin{proof}
We have for any $e \in E_n(R)$ and $g \in GL_n(R)$,
\[
\pmat {geg^{-1}} 0 0 1 = \pmat g 0 0 {g^{-1}} \pmat e 0 0 1 \pmat {g^{-1}} 0 0 g 
\]
Since each of the elements are in $E_{2n}(R)$, the product is in $E_{2n}(R)$.
This implies $E(R)$ is normal in $GL(R)$.
\end{proof}

Now we arrive at our definition of $K_1(R)$.
\begin{definition}[$K_1(R)$]
 We define group $K_1(R)$ to be the quotient
 \[
 K_1(R) := GL(R)/E(R) 
 \]
\end{definition}

\begin{example}
If $F$ is a field, then $K_1(F) = F^\times$.
\end{example}

\subsection{Properties of $K_1$}

\begin{proposition}\label{prop:element}
For elementary matrices, we have $e_{ij}(r) e_{ij}(s) = e_{ij}(r+s)$, and $e_{ij}(r)$ has the inverse $e_{ij}(-r)$.
Otherwise, we have the following commutation relations between arbitrary elementary matrices:
\[
[e_{ij}(r),e_{kl}(s)] = \begin{cases}
1 \quad &\text{if $j \neq k$ and $i \neq l$} \\
e_{il}(rs) \quad &\text{if $j = k$ and $i \neq l$} \\
e_{kj}(-sr) \quad& \text{if $j \neq k$ and $i = l$}
\end{cases}
\]
\end{proposition}

We say that a group $G$ is \emph{perfect} if $G = [G,G]$. Proposition \ref{prop:element}
immediately gives us that
\begin{lemma}
If $n \geq 3$, then $E_n(R)$ is perfect.
\end{lemma}
\begin{proof}
If $i,j,k$ are distinct, we have $e_{ij}(r) = [e_{ik}(r),e_{kj}(1)]$ which shows $E_n(R) \subseteq [E_n(R),E_n(R)]$,
giving us the result.
\end{proof}

This also shows that the commutator subgroup $[GL(R),GL(R)]$ contains all elements of $E(R)$
after stabilizing, since every element in $E(R)$ can be thought of as an element in $E_n(R)$ for some large $n$.
Here is a result due to Whitehead's 1950 paper [Whitehead].

\begin{lemma}[Whitehead's Lemma]
We have $E(R) = [GL(R), GL(R)]$. In particular, $E(R)$ is perfect, and $K_1 = GL(R)_{ab}$.
\end{lemma}
\begin{proof}
Since the commutator subgroup contains $E(R)$, it suffices to show the converse. 
Every commutator in $GL_n(R)$ may be written as the following product
\[
[g,h] = \pmat g 0 0 {g^{-1}} \pmat h 0 0 {h^{-1}} \pmat {{(hg)}^{-1}} 0 0 {hg}
\]
But each of the elements in the product are in $E_{2n}(R)$.
This gives us the desired claim.
\end{proof}

Clearly, this also means that $K_1(R)$ is abelian for any ring $R$. 

\subsection{Characterization of $K_1$ in terms of automorphisms}

\begin{definition}
Define the category $\mathbf P(R)$
\end{definition}

\begin{theorem}
$K_1(R)$ may also be described as the group generated by symbols 
$[(P,\alpha)]$ where $P$ is a (right) projective module over $R$ and $\alpha : P \xrightarrow{\cong} P$
is an automorphism, subject to the following relations:
\begin{enumerate}
  \item $[(P,\alpha)] = [(Q,\beta)]$ if there is an $f : P \xlongrightarrow{\cong} Q$, 
  for which 
  $f \alpha = \beta f$.
  \item $[(P,\alpha \circ \beta)] = [(P,\alpha)] [(P,\beta)]$
  \item $[(P \oplus Q, \alpha \oplus \beta)] = [(P,\alpha)] [(Q,\beta)]$
\end{enumerate}
\end{theorem}

Note that the relations force $K_1(R)$ to be abelian.

\subsection{Homological Interpretaion of $K_1$}
%The group $K_1$ was first studied by Whitehead in the context of homotopy torsion.

$K_1$ may be interepreted in the language of homological algebra. We refer the reader to [Weibel, HomoAlg]
Note that the first group homology $H_1(G;\Z)$ of any group $G$
is naturally isomorphic to $G/[G,G]$. This gives us
\[
K_1(R) = H_1(GL(R);\Z)
\]

\begin{theorem}[Bass]
We have $K_1(R) \cong \varinjlim\limits_{P \in t \mathbf P} H_1(\Aut(P);\Z)$
\end{theorem}

\begin{remark}
The group \(K_1\) was studied by Whitehead in the theory of simple homotopy types and whitehead torsion.
The Whitehead group, which arises naturally in this context, is a quotient of \(K_1(\mathbb Z[\pi])\).
The Whitehead group $Wh_1(\pi)$ of a (fundamental) group $\pi$ is the quotient $K_1(\Z[\pi])/H$ where $H$ is the subgroup
of $K_1(\Z[\pi])$ generated by $+- 1$ and $GL_1(\pi)$.
\end{remark}



%\subsection{Fundamental Thoerems for $K_0$ and $K_1$}

%\clearpage

\section{$K_2$ of a ring}

The notion of the $K_2$ group of a ring was introduced by Milnor in 1967 [].
We refer the reader to [Milnor Book] for a detailed discussion of the subject.
Here, we shall define $K_2$ and state important results.

We showed that the group $E_n(R)$ obeyed certain relations in \ref{prop:element}

\begin{definition}[$K_2(R)$]
Let $n \geq 3$.
The \emph{Steinberg group} $St_n(R)$ of a ring $R$ is defined as the group generated
by the symbols $x_{ij}(r)$ for $1 \leq i \neq j \leq n$ and $r \in R$
subject to the relations
\begin{align}
  x_{ij}(r) x_{ij}(s) = x_{ij} (r + s) 
\end{align}
\begin{align}
  [x_{ij}(r),x_{kl}(s)] = \begin{cases}
  1 \quad &\text{if $j \neq k$ and $i \neq l$} \\
  x_{il}(rs) \quad &\text{if $j = k$ and $i \neq l$} \\
  x_{kj}(-sr) \quad& \text{if $j \neq k$ and $i = l$}
  \end{cases}
\end{align}
\end{definition}

Note that these relations are satisfied by elementary matrices.
This gives us a surjective homomorphism $St_n(R) \to E_n(R)$.

It is clear we have another filtered system given by embeddings $St_n(R) \hookrightarrow St_{n+1}(R)$,
the colimit of which we call the \emph{stable Steinberg group}, denoted $St(R)$.

Since we can think of $E_n(R)$ as living in $E(R)$, we have homomorphisms $St_n(R) \to E_n(R) \to E(R)$,
which induce a (surjective) group homomorphism $\pi : St(R) \to E(R)$.

\begin{definition}[$K_2$ of a ring]
We define the group $K_2(R)$ for a unital ring $R$ as the kernel of $\pi : St(R) \to E(R)$.
\end{definition}
Note that this gives us an exact sequence
\begin{equation}
1 \to K_2(R) \to St(R) \xrightarrow{\pi} E(R) \to 1 
  \label{eq:K2exact}
\end{equation}
which completes into a longer sequence
\begin{equation}
1 \to K_2(R) \to St(R) \to GL(R) \to K_1(R) \to 1 
  \label{eq:K2longerexact}
\end{equation}

We state a result due to Steinberg. (See [])
\begin{theorem}[Steinberg]
$K_2(R)$ is precisely the center of $St(R)$, hence abelian.
\end{theorem}

\subsection{Universal Central Extensions}

The Steinberg group $St(R)$ can be described in terms of universal 
central extensions, which we shall define. See [Milnor] and [Suz].
From the viewpoint of homological algebra, central extensions of $G$ are classified by
the cohomology group $H^2(G;\Z)$ [See WHomo], and it turns out one can describe $K_2(R)$
as a homology group, $K_2(R) \cong H_2(E(R);\Z)$.
Here, we give a skeleton of the results leading up to this fact.

\begin{definition}
 A \emph{central extension} of a group $G$ is an exact sequence
 \[
 \grpexact A X G
 \]
 such that the image of $A$ is in the center of $H$.
\end{definition}

We will often simply write $X \xrightarrow{\pi} G$ to denote the central extension when $A$ is clear.

\begin{definition}
A \emph{univeresal central extension} of $G$ is a central extension which is universal among central extensions of $G$. That is,
for any central extension $Y \xrightarrow{\tau} G$ we have a unique group homomorphism $f : X \to Y$
\[
\begin{tikzcd}[row sep=large, column sep=large]
X \arrow[r, "\pi"] \arrow[d, dashed, "\exists!f"'] & G \\
Y \arrow[ur, "\tau"']
\end{tikzcd}
\]
\end{definition}
Clearly, a universal central extension is unique up to a unique isomorphism if it exists.

\begin{lemma}
If $G$ has a universal central extension $X \xrightarrow{\pi} G$, then $G$ and $X$ are both perfect.
\end{lemma}
\begin{proof}
If not, then $X/[X,X]$ is nontrivial, and we have 
\end{proof}

We also have a converse. 

Let a group $G$ have a presentation $F/R$ where $F$ is a free group and $R$ is the subgroup of relations.
This may be described as an exact sequence
\[
\grpexact R F G
\]
Since $[R,F]$ is a normal subgroup of $F$, we get the following central extensions.
\[
  \grpexact {R/[R,F]} {F/[R,F]} {G} \\
\]
\begin{equation}
  \grpexact {(R \cap [F,F])/[R,F]} {[F,F]/[R,F]} {[G,G]}
  \label{eq:beforeH2}
\end{equation}

Note the group $(R \cap [F,F])/[R,F]$ is exactly the homology group $H_2(G;\Z)$.

\begin{theorem}[Recognition Theorem for Universal Central Extensions]\label{thm:recog_universalcentral}
Every perfect group $G$ has a universal central extension
\[
1 \to H_2(G;\Z) \to [F,F]/[R,F] \to G \to 1
\]
Consequently, if $X$ is any central extension of $G$, then the following are equivalent
\begin{enumerate}
  \item $X$ is a universal central extension.
  \item $X$ is perfect and every central extension of $X$ splits.
  \item $H_1(X;\Z) = H_2(X;\Z) = 0$
\end{enumerate}
\end{theorem}
See [Milnor] for a complete proof.
The following result follows. (See [Kervaire, Steinberg])
\begin{theorem}[Kervaire, Steinberg]
The Steinberg group $St(R)$ is the universal central extension of $E(R)$.
Thus we can conclude
\[
  K_2(R) \cong H_2(E(R);\Z)
\]
\end{theorem}

\subsection{Bass's $K_2$}
We also state the definition of $K_2(R)$ due to Bass,
which turns out to be equivalent to our definition of $K_2(R)$. (See [Bass], [Kervaire-Steinberg])
\[
K_2(R) \cong \varinjlim\limits_{P \in t\mathbf P} H_2([\Aut(P),\Aut(P)];\Z)
\]

In the next chapter, we shall discover homotopy-theoretic constructions
which yield the groups $K_n(R)$ and prove their equivalence with these classical notions.

